\appendix
\section*{Appendix E — Reversible Quest Algebra}
\addcontentsline{toc}{section}{Appendix E — Reversible Quest Algebra}

\section{Quest Path Framework}

\subsection{Operator Paths}
A reversible quest is represented as an operator path
\[
\gamma = (O_1, O_2, \ldots, O_n),
\]
where each $O_i$ is a reversible operator acting on the state space
$\mathcal{S}$.

\subsection{Forward Traversal}
Forward traversal applies operators in sequence:
\[
\gamma(s) = O_n(\cdots O_2(O_1(s))\cdots).
\]

\subsection{Reverse Traversal}
Reverse traversal applies inverses in reverse order:
\[
\gamma^{-1}(s)
    = O_1^{-1}(O_2^{-1}(\cdots O_n^{-1}(s)\cdots)).
\]

\section{Reversibility and Geometry}

\subsection{State Reversibility}
For any state $s$,
\[
\gamma^{-1}(\gamma(s)) = s.
\]
State reversibility is guaranteed by operator inverses.

\subsection{Geometric Irreversibility}
Curvature and frame shifts do not reverse:
\[
F' = \mathrm{Hol}(\gamma)(F),
\qquad
F'' = \mathrm{Hol}(\gamma^{-1})(F'),
\]
and in general $F'' \neq F$.

\section{Curvature Accumulation}

\subsection{Holonomy Contributions}
Curvature arises from operator non–commutativity:
\[
H_{ij} = [O_i, O_j]
       = O_j \circ O_i - O_i \circ O_j.
\]

\subsection{Path Curvature}
Total curvature accumulated along a quest path is
\[
K(\gamma) = \sum_{i < j} H_{ij}.
\]

\subsection{Quest Energy}
Define quest energy as
\[
E(\gamma) = \|K(\gamma)\|.
\]

\section{Definitions}

\begin{definition}[Reversible Quest]
A reversible quest is an operator path whose forward and reverse
traversals restore the original state.
\end{definition}

\begin{definition}[Quest Curvature]
Quest curvature is the accumulated holonomy contribution
\[
K(\gamma) = \sum_{i < j} H_{ij}.
\]
\end{definition}

\begin{definition}[Quest Energy]
Quest energy is the norm of the curvature tensor:
\[
E(\gamma) = \|K(\gamma)\|.
\]
\end{definition}

\begin{definition}[Quest Composition]
Two quests compose as
\[
\gamma_2 \circ \gamma_1
    = (O_1, \ldots, O_m, O_{m+1}, \ldots, O_n).
\]
\end{definition}

\section{Theorems}

\begin{theorem}[State Reversibility]
For any quest path $\gamma$ and state $s$,
\[
\gamma^{-1}(\gamma(s)) = s.
\]
\end{theorem}

\begin{proof}
Each operator $O_i$ has an inverse $O_i^{-1}$ satisfying
$O_i^{-1}(O_i(s)) = s$. Applying inverses in reverse order restores the
original state.
\end{proof}

\begin{theorem}[Curvature Non-Reversibility]
Curvature does not reverse under path inversion:
\[
K(\gamma^{-1}) \neq -K(\gamma).
\]
\end{theorem}

\begin{proof}
Curvature depends on commutators $[O_i, O_j]$. Reversing the path changes
operator order but does not negate commutators. Thus curvature accumulates
independently of state reversibility.
\end{proof}

\begin{theorem}[Additivity of Curvature]
For paths $\gamma_1$ and $\gamma_2$,
\[
K(\gamma_2 \circ \gamma_1)
    = K(\gamma_1) + K(\gamma_2).
\]
\end{theorem}

\begin{proof}
Curvature is defined as a sum over operator pairs. Concatenating paths
adds new operator pairs, producing additive curvature.
\end{proof}

\section{Interaction with Spectral Ecology}

\subsection{Curvature-Induced Laplacian Warp}
Curvature modifies adjacency weights:
\[
w_{ij}' = w_{ij} + H_{ij},
\]
producing a warped Laplacian
\[
L' = D' - W'.
\]

\subsection{Spectral Drift}
Warping the Laplacian produces spectral drift:
\[
(\lambda_k, u_k) \mapsto (\lambda_k', u_k').
\]

\subsection{Ecological Consequences}
Quest curvature induces:
\begin{itemize}
    \item monster-family evolution,
    \item faction-biased ecological zones,
    \item narrative-tension propagation,
    \item curvature-sensitive rare species.
\end{itemize}

\section{Discussion}
This appendix formalizes reversible quest algebra, including operator
paths, curvature accumulation, geometric irreversibility, and spectral
interaction. These structures integrate with holonomy curvature,
Laplacian geometry, and spectral ecology in the main preprint.
